\documentclass[a4paper,12pt]{article}
\usepackage[utf8]{inputenc}
\usepackage[T1]{fontenc}
\usepackage[ngerman]{babel}
\usepackage{amsmath}
\usepackage{amssymb}
\usepackage{enumitem}
\usepackage{geometry}
\geometry{a4paper, margin=2.5cm}

\title{Einsendeaufgaben -- Aufgabe 5.2\\[0.5em]
\large Modul 61111: Mathematische Grundlagen}
\author{}
\date{}

\begin{document}
\maketitle

\section*{Aufgabe 5.2}

\begin{enumerate}[label=\alph*)]

    \item
    Die Folge $(x_n)$ mit $x_n = \frac{1}{n}$ konvergiert gegen $0$.

    Da für jedes $x_n$ gilt $x_n \neq 0$, folgt
    \[
        \lim_{n \to \infty} f(x_n) = \lim_{n \to \infty} 1 = 1 \neq 0 = f(0).
    \]

    Also ist $f$ in $0$ nicht stetig.

    \item
    Sei $\varepsilon > 0$ beliebig. Dann wählen wir
    ein $\delta := \varepsilon^2$.

    Sei $x \in \mathbb{R}$ mit $|x| < \delta$ beliebig.

    Dann gilt:
    \begin{align*}
        |f(x) - f(0)| = |\sqrt{x}| = \sqrt{x}
        < \sqrt{\varepsilon^2} = \varepsilon.
    \end{align*}

    Nach dem $\varepsilon$-$\delta$-Kriterium ist $f$ in $0$ stetig.

\end{enumerate}

\end{document}
